\documentclass[main.tex]{subfiles}

\begin{document}

\section{Theoretical Framework}
\label{sec:redistribution}

In this section, we describe two complementary approaches to quantify the consumer welfare consequences of interchange fees: a sufficient statistic approach and a structural model.
Both approaches capture how consumer sorting and merchant fee heterogeneity shape the redistributive effects of interchange fees.
They also allow us to evaluate different regulatory counterfactuals such as European-style fee regulation or the Durbin Amendment.

Each approach rests on different assumptions. 
The sufficient statistic approach applies broadly across utility functions and market structures, and highlights the first-order determinants of the consumer welfare effects of interchange fees.
The main advantage of the sufficient statistic approach is that it avoids strong functional-form assumptions, but without additional structure governing merchant competition, we must make an assumption about how interchange fees pass through to retail prices. 
The structural model imposes more parametric assumptions on consumer utility functions and the nature of competition between firms. 
This allows us to relax the assumption of full pass-through and capture how large merchants pass through interchange fees imperfectly to retail prices due to market power.
Because the estimates from the two methods are quantitatively similar, we present the more intuitive sufficient statistic approach first.

\subsection{A Sufficient Statistic Approach}

We model the behavior of consumers, merchants, and payment networks.
Consumers with fixed payment preferences choose how to allocate spending across merchants subject to a budget constraint.
Merchants set retail prices in response to the composition of payments they receive and the fees on each payment method.
Payment networks set interchange fees and pay consumers rewards.

\subsubsection{Consumers}
The sufficient statistics economic environment places few restrictions on consumer preferences.
Consumers with fixed payment preferences allocate spending across merchants.
Let $k=1,\ldots,K$ denote different payment methods, and let the share of consumers who use payment method $k$ equal $\mu_k$, with $\sum_k \mu_k=1$.
Type-$k$ consumers have income $1 + f_k$, where $f_k$ denotes rewards.
Consumers spend their income on goods from $j=1,\ldots,J$ merchants.
Let $p_j$ denote the price set by merchant $j$ and $q_{jk}$ denote per-capita purchases from merchant $j$ by payment-method-$k$ consumers.
We use stars to denote equilibrium values.

Consumers solve the utility maximization problem

\[q^{*}_{jk} = \arg\max_{q_{jk}\geq 0} U_k\left(\{q_{jk}\}_{j}\right) \text{ s.t. } \sum_j p_j q_{jk} \leq 1 + f_k,\]
where $U_k$ is a strictly increasing, strictly concave, and twice-continuously differentiable utility function.\footnote{ These assumptions rule out perfect substitutes and ensure that demand functions are continuously differentiable almost everywhere.
We assume own-price elasticities dominate cross-price elasticities (Assumption \ref{assump:dominant-diagonal} in Appendix \ref{sec:appendix-theory}) and are not too large relative to the inverse of the maximum merchant fee. } Allowing utility functions to depend on payment preferences enables the model to capture that credit, debit, and cash consumers may shop at different stores.

\subsubsection{Merchants}

Merchants set a uniform price for all consumers, regardless of payment method, consistent with the institutional reality that most U.S. merchants do not price discriminate by payment type \citep{Stavins2018}.\footnote{Individual merchants face weak incentives to surcharge by payment method despite different acceptance costs.
Customer backlash creates first-order costs \citep{Stavins2018,Caddy.etal2020}, while a standard envelope theorem argument shows gains from price discrimination are second-order in merchant fees \citep{Wang2025}.}
Let $\tau_{jk}$ denote the interchange fee charged to merchant $j$ for card $k$, expressed as a fraction of the transaction value. For empirical and expositional convenience, we assume that interchange fees are a function of merchant sector and size. As described in Sections \ref{sec:institutional-data} and \ref{sec:costs}, this assumption is largely consistent with the institutional setting and data.
Let $s_j$ denote the sector of merchant $j$, and let $b_j = 1$ if the merchant has a negotiated discount. The interchange fee charged to merchant $j$ for payment method $k$ is calculated as follows:
\begin{equation}
\tau_{jk}=\overline{\tau}_{k}+\tau_{s_{j}k}+b_{j}\delta_{k}\label{eq:interchange-fee-model}
\end{equation}
where $\overline{\tau}_{k}$ is the baseline fee for payment method $k$, $\tau_{s_{j}k}$ is a sector adjustment, and $\delta_{k}<0$ is a negotiated discount available to large merchants.\footnote{This specification abstracts from the role of ticket size in shaping average transaction fees, as we do not separately model regulatory shocks to the fixed versus linear components of interchange fees.}

Rather than explicitly modeling merchant competition or imposing functional forms on the demand side, we assume that prices are log-linear in weighted interchange fees, and that merchants pass through interchange fees at a rate $\rho^P\in[0,1]$.
\begin{equation}
\log p_{j}=\log\overline{p}_{j}+\rho^P\sum_{k}\mu_{jk}\tau_{jk},\label{eq:retail-price-pass}
\end{equation}
where $\overline{p}_j$ is a baseline price in the absence of interchange fees, and $\mu_{jk}$ is the share of merchant $j$'s sales paid with method $k$. Equivalently, $\mu_{jk}=\mu_k q_{jk}^{*}/Q_{j}^{*}$ where $Q_{j}^{*}=\sum_{k}\mu_k q_{jk}^{*}$ is total sales at merchant $j$.
Under monopolistic competition with CES preferences, this equation is a first-order approximation for the true pricing function with $\rho^P=1$.

\subsubsection{Card Networks and Rewards}
We assume that card networks pass interchange fees through to lump-sum consumer rewards at a constant \emph{marginal} rate $\rho^R\in[0,1]$, on top of a baseline reward $\bar f_k$.\footnote{Lump-sum rewards and merchant-specific or sector-specific rewards have identical first-order consumer welfare effects.
By Roy's identity, the welfare impact of reallocating spending across merchants in response to category-based rewards is second-order; the first-order effect is the income change. }
Rewards are then determined by:
\begin{align}
f_{k} & =\bar f_{k}+\rho^R\,\bar\tau^{k},
\qquad
\bar\tau^{k}\equiv\frac{\sum_{j}p_{j}q_{jk}^{*}\,\tau_{jk}}{\sum_{j'}p_{j'}q_{j'k}^{*}}.
\label{eq:reward-pass-through}
\end{align}
Here $\bar\tau^{k}$ is the spending-weighted average interchange fee that method-$k$ users encounter, and $\bar f_k$ is a baseline reward that absorbs every component that does not move with the fee: network processing costs, issuer margins from market power, and cross-subsidies from other bank products such as interest income on revolving balances \citep{Agarwal.etal2025, Agarwal.etal2026} or cross-sold deposits \citep{Brown.McDonald2026}.
Because income is normalized to one, $f_k$ equivalently represents the per-dollar reward rate or a lump-sum transfer to type-$k$ consumers.

We interpret rewards broadly to include any subsidy for card use.
For example, \citet{Kay.etal2018, Mukharlyamov.Sarin2025} show that high pre-Durbin debit interchange fees helped to fund free checking accounts.
To the extent that certain credit cards charge annual fees, these rewards are net of those fees.

\subsubsection{Sufficient Statistics for Consumer Welfare Effects}

We develop sufficient statistics to conduct an equilibrium welfare evaluation of interchange fees across consumer groups when consumer payment choices and merchant acceptance decisions are held fixed, but prices and consumption can re-adjust, isolating redistribution across groups holding the payment mix fixed.
Section \ref{subsubsec:robustness} discusses the robustness of this assumption; Appendix \ref{sec:appendix-endogenous-choice} allows payment choice to respond endogenously to rewards.
The key insight is that first-order consumer welfare effects depend on the variation in the composition and cost of payments across merchants, rather than on demand elasticities. We use this framework to calculate the welfare and redistributive effects of eliminating interchange fees in Section \ref{sec:Results}.

\begin{theorem}[Sufficient Statistic for Consumer Welfare Redistribution]\label{thm:sufficient-stat}
Let $x_l$ be any parameter that shifts $\tau_{jl}$ (e.g., the baseline fee, sector adjustment, or negotiated discount), and let $\tau_{\max}=\max_{j,k}|\tau_{jk}|$ denote the largest interchange fee in absolute value across all merchants and payment methods.
There is a threshold $\bar\tau>0$, depending only on demand primitives, such that whenever $\tau_{\max}<\bar\tau$ the first-order money-metric welfare effect for consumers of type $k$ is
\begin{equation}
\frac{1}{\lambda_k}\frac{dV_k}{dx_l}
= \mu_k^{-1}\!\left(
-\,\rho^P \mathbb{E}_R\!\left[\mu_{jk}\mu_{jl}\tfrac{\partial\tau_{jl}}{\partial x_l}\right]
+ \mathbbm{1}\{k=l\}\,\rho^R \mathbb{E}_R\!\left[\mu_{jk}\tfrac{\partial\tau_{jk}}{\partial x_k}\right]
\right) + O(\tau_{\max}),\label{eq:sufficient-stat-formula}
\end{equation}
where $\lambda_k$ is the marginal utility of income and $\mathbb{E}_R\!\left[\mu_{jk}\right]$ denotes expectations weighted by each merchant's share of revenue $\omega_j = R_j / \sum_{j'} R_{j'}$ in the overall economy.
The total effect of a discrete fee component can be calculated by multiplying by $x_l$.
Multiplying by $\mu_k$ yields the aggregate dollar welfare effect for group $k$.
The threshold $\bar\tau$ is characterized explicitly in the proof (Appendix~\ref{sec:appendix-theory}, Step~C).

\end{theorem}

\subsubsection{Intuition and proof sketch.}

The theorem states that two channels determine how interchange fees redistribute welfare when consumer payment choices and merchant acceptance decisions are fixed.
The first term is the retail-price channel: When interchange fees on payment method $l$ increase, retail prices rise at merchants where method-$l$ users shop.
The size of the mechanical price change, holding fixed the composition of payments, is precisely $\mu_{jl} \frac{\partial \tau_{jl}}{\partial x_l}$.
Consumers using payment method $k$ are harmed in proportion to their overlap with method-$l$ users, which then yields the product $\mu_{jk}\mu_{jl} \frac{\partial \tau_{jl}}{\partial x_l}$.
The second term is the direct reward channel: method-$l$ users receive more in card rewards when their interchange fees rise, affecting only those users.

The sufficient statistic captures two key insights about the effects of interchange fees on consumer welfare.
First, it shows that consumer substitution across merchants does not matter to first-order.
Even though the statistic is expressed as a mechanical calculation holding payment composition fixed, it captures the full first-order welfare effects even as consumers adjust their consumption decisions in response to changes in retail prices.
Second, it shows that the effects of interchange fees on consumers who use different payment methods can be measured with merchant-level payment shares; no consumer-level data is needed.

The main challenge in proving the theorem is showing that the equilibrium changes in retail prices across all merchants are well approximated by the mechanical effects that fix consumption decisions.
In our setting, merchants' prices depend on the composition of consumers at each merchant.
As prices change, compositions change as well.
While Roy's identity tells us that the changes in retail prices are sufficient to characterize the welfare effects of interchange fees, it provides no guarantee that the mechanical effects described above approximate the final change in retail prices.
We show that as long as price elasticities and income semi-elasticities and are not too large relative to the level of fees, then these additional feedback effects are second-order.

\subsubsection{Sufficient statistic parameters.}

To calculate the sufficient statistics for consumer welfare as per eq. (\ref{eq:sufficient-stat-formula}), we need to impose assumptions about the pass-through of interchange fees to retail prices and rewards, measure the revenue-weighted first and second-moments of payment shares across merchants, and estimate the fee parameters.
In our baseline implementation, we assume full pass-through of interchange fees to retail prices, so $\rho^P=1$, and full pass-through of interchange fees to rewards, so $\rho^R=1$. 
Section \ref{subsubsec:robustness} discusses the robustness of our results to these assumptions.
We recover the revenue-weighted first and second-moments with their empirical analogues:
\[
\widehat{\mu}_k = \sum_j \widehat \omega_j\,\mu_{jk}, \qquad
\widehat{\mathbb{E}_R[\mu_{jk}\mu_{jl}]} = \sum_j \widehat \omega_j\,\mu_{jk}\mu_{jl}.
\]
Both sets of moments are straightforward to compute using our merchant-level data.
We recover the fee parameters by regressing firm-level interchange fees for each payment method on sector dummies and a dummy for having more than \$1 billion in sales in the grocery and retail sectors.
We focus on large firms in these sectors, where discounts are most pronounced in the data (Section \ref{sec:costs}).
The regression constant yields the baseline fee $\overline{\tau}_k$; the sector dummy coefficients yield the sector adjustments $\tau_{sk}$; and the bargaining dummy coefficient yields the negotiated discount $\delta_k$.

\subsubsection{Relation to alternative approaches.}

Our sufficient-statistic approach yields a simple formula that requires only the payment shares $\mu_k$, the overlap moments $\mathbb{E}_R[\mu_{jk}\mu_{jl}]$, and the fee parameters.
In contrast, a fully structural demand estimation requires specifying and estimating utility functions, demand elasticities, and merchant pricing behavior across all consumer types and merchants.
The theorem implies that consumer substitution patterns do not affect first-order redistribution.
To validate this sufficient statistic approach, we next calibrate a structural model that matches key data moments while allowing for imperfect pass-through.

\subsection{Structural Model}

We calibrate a structural model of retail competition to validate the sufficient statistic approach.
While the sufficient-statistics approach provides intuitive and transparent estimates, there are two potential sources of error..
First, the first-order approximation may introduce a meaningful error by omitting second-order utility losses from consumer reallocation.
Second, the assumptions about the pass-through of interchange fees to retail prices rule out realistic forces such as strategic complementarities in merchant pricing.
By calibrating a structural model, we can decompose potential sources of error into these two sources.
The models are calibrated to match key moments from the Fiserv data on payment composition and costs across merchants and sectors, ensuring that the model reproduces the heterogeneity central to our empirical findings.

\subsubsection{Consumer Preferences}

We model consumer demand using a nested CES structure that allows for rich heterogeneity in shopping patterns across payment methods.
Consumers with different payment methods have different tastes for merchants, and this heterogeneity matches the observed variation in the composition of payments across merchants, linking the model directly to the empirical dispersion documented in the data.

In a representative market $m$, consumers who use payment method $k=1,\dots,K$ have Cobb-Douglas preferences over sectors $s=1,\dots,S$:
\[
\log U_{k}=\sum_{s=1}^{S}\alpha_{ks}\log U_{ks},\quad\sum_{s=1}^{S}\alpha_{ks}=1
\]
Utility in a sector is a CES aggregator over the $J$ firms serving that market:
\[
U_{ks}^{\frac{\sigma-1}{\sigma}}=\sum_{j=1}^{J}a_{jks}^{\frac{1}{\sigma}}q_{jks}^{\frac{\sigma-1}{\sigma}}
\]
where the $a_{jks}$ are taste shifters that allow us to match the distribution of payment shares and firm size. Consumers solve an optimal consumption problem, in which their budget constraint reflects their rewards:
\begin{equation}
\max_{q_{k}}U_{k}\quad\text{s.t. }\sum_{s}\sum_{j}q_{jks}p_{js}\leq1+f_{k} \label{eq:consumption-problem}
\end{equation}

Intuitively, a sector with a high $\alpha_{ks}$ has a high utility for consumers $k$.
For example, premium card consumers may have a higher utility for leisure travel than cash consumers.
Similarly, within a particular sector, consumers' preferences over firms are correlated with their payment methods, which is reflected in $a_{jks}$.
For example, premium card consumers may prefer to shop for groceries at Whole Foods relative to a discount grocer, while cash consumers may prefer the opposite (given prices).
This rich heterogeneity allows firms to have different market shares across consumers who use different payment methods.
The parameter $\sigma$ captures consumers' willingness to substitute across goods (i.e., consumers' price elasticity), so a higher $\sigma$ implies more price-elastic consumers. 

The standard solution yields:
\begin{align}
q_{jks} & =\alpha_{ks}\times\left(1+f_{k}\right)\times\frac{a_{jks}p_{js}^{-\sigma}}{P_{ks}^{1-\sigma}} \\ \label{eq:consumer-demand}
P_{ks}^{1-\sigma} & =\sum_{j=1}^{J}a_{jks}p_{js}^{1-\sigma}
\end{align}
The expression is standard and intuitive.
For a given price, goods with higher $a$ and sectors with higher $\alpha$ have higher demand.
Higher relative prices lower demand, and higher price elasticity ($\sigma$) results in larger declines in demand given a price increase.

\subsubsection{Merchant Problem}

A mix of large and small merchants maximize profits by setting a uniform price for all consumers.
A critical feature of our model is that it allows a discrete number of firms in a representative market, following \citet{Hottman.etal2016}, which generates both imperfect pass-through of idiosyncratic cost shocks and strategic complementarities in pricing.
By varying the number of large firms per market, we can control the degree of pass-through and thereby assess the sensitivity of our redistribution results to this margin, linking model structure directly to the empirical pass-through question.

Formally, let the measure of consumers who use payment method $k$ equal $\mu_k$.
For convenience, we normalize marginal costs to 1, as CES models do not separately identify costs and taste shifters and we do not observe cost data,
Then merchants solve:
\[
p_{js}=\argmax_{p}\sum_{k}\mu_k q_{jks}\left(p\left(1-\tau_{jks}\right)-1\right)
\]

Because firms are discrete, they internalize the effect of their own price setting.
Large firms thus exert market power, charge higher markups, and feature incomplete pass-through of interchange fees.
Define $\rho_{jks}$ to be the payment-method-$k$ revenue share of firm $j$ among all firms in that sector:
\[
\rho_{jks} =\frac{q_{jks}p_{js}}{\sum_{i}q_{iks}p_{is}}
\]
The optimal price then equals
\begin{align}
  \label{eq:optimal-price-structural}
p_{js} & = \frac{\sum_{k}\mu_{jks}\varepsilon_{jks}}{\sum_{k}\mu_{jks}\left(\varepsilon_{jks}-1\right)\left(1-\tau_{jks}\right)} \\
  \varepsilon_{jks} & = \sigma(1-\rho_{jks}) + \rho_{jks} \\
\mu_{jks} & =\frac{\mu_k q_{jks}}{\sum_{l}\mu_l q_{jls}}
\end{align}

In this model, markups are increasing in firm size and the pass-through of idiosyncratic cost shocks is decreasing in firm size. 
When consumers have CES preferences and firms are monopolistically competitive, all merchants fully pass through idiosyncratic cost shocks while ignoring the pricing decisions of their competitors. 
However, if firms are large, they internalize the effect of their pricing decisions on the aggregate price index; they only partially pass through idiosyncratic cost shocks. 
Their prices also depend on the prices that other merchants set, i.e., they exhibit strategic complementarity \citep{Amiti.etal2019}.

\subsubsection{Network Problem}

We model network price setting in a mechanical manner.
As in the sufficient statistic framework, card networks rebate interchange fees to consumers as rewards, net of processing costs and a fixed arithmetic margin.
Networks mechanically pass on merchant fees net of costs in the form of rewards:
\[
f_{k}=\frac{\sum_{s}\sum_{j}q_{jks}p_{js}\left(\tau_{jks}-c_{k} - m_k \right)}{\sum_{s}\sum_{j}q_{jks}p_{js}}
\]
Here, we focus on the case of full pass-through as the main purpose of the structural model is to evaluate the effects of ignoring consumer reallocation across merchants and potential strategic complementarities in merchant pricing, not to model network conduct.

Merchant fees at each firm are based on card type, sector, and firm size. The interchange fee schedule mirrors the empirical structure documented in \ref{sec:costs}: fees vary by card type (regulated debit, exempt debit, basic credit, premium credit), by merchant sector, and by merchant size (with discounts for large firms).
The firm size discount is implemented by an indicator $\delta_{j}$ for whether the firm has more than $1$ billion in sales and a card-specific discount $b_{k}<0$:
\[
\tau_{jks}=\tau_{k}+\tau_{ks}+\delta_{j}b_{k}
\]

\subsubsection{Equilibrium}

An equilibrium is defined by a vector of prices $p_{js}$ for each sector and a vector of consumption choices $q_{jks}$ for consumers of each payment method, such that consumers maximize utility given prices as per Equation \ref{eq:consumption-problem} and merchants maximize profits as per Equation \ref{eq:optimal-price-structural}.

\subsubsection{Calibration}

We next calibrate our model to the data.
We target moments describing the composition of payments across merchants and sectors.
We calibrate most parameters sector-by-sector to allow for the substantial cross-sector heterogeneity documented in Section \ref{sec:costs}.
First, we calibrate the $\alpha_{ks}$ parameters to match the allocation of spending across sectors.
We then parameterize the taste shifters $\log a_{jks}$ as jointly normal with a correlation structure to match the revenue-weighted mean and covariance matrices of payment shares across merchants in each sector.
For the grocery and merchandise sectors, we allow for the distributions of $a_{jks}$ to differ for small and large firms (who receive negotiated interchange).
We use the same estimated fee parameters as in the sufficient statistics exercise.
Appendix \ref{sec:appendix-calibrated} provides full calibration details and goodness-of-fit figures, showing that the model closely matches the empirical moments.

\end{document}